\chapter{Introduction}
\label{ch:introduction}
Large Language Models have seen rising adoption in recent years, even in high-stakes environments \cite{maityLargeLanguageModels2025, laiLargeLanguageModels2024, hungWalkingTightropeEvaluating2023}, but their opaque nature remains challenging: in most cases, it is not possible to assess how and why a model arrived at a given prediction \cite{qamarUnderstandingBlackboxInterpretable2023}. Consequently, there is an ongoing research effort to make models explainable, to render their inner workings transparent, to build trust, and, in some cases, to fulfill regulatory requirements \cite{hsiehComprehensiveGuideExplainable2024, aliExplainableArtificialIntelligence2023, herreraMakingSenseUnsensible2025}.

Some approaches to \emph{explainability} seek to visualize the path a decision has taken throughout the model using technical means, e.g., by showing how certain attention weights or neuron activations influenced the result \cite{nazirSurveyFeatureAttribution2025, zhaoExplainabilityLargeLanguage2024}. However, not only does the implementation of such techniques depend on the model's architecture, but they fundamentally require the model to be open to such analysis \cite{nazirSurveyFeatureAttribution2025}. Other techniques, such as \emph{LIME} \cite{ribeiroWhyShouldTrust2016} and \emph{SHAP} \cite{lundbergUnifiedApproachInterpreting2017}, instead measure \emph{Feature Attribution}, i.e., how strongly individual parts of the input influence the output. While this approach is model-agnostic, it comes with its own set of problems, such as high computational overhead and considerable instability with respect to small input changes \cite{nazirSurveyFeatureAttribution2025}.

One major advantage of LLMs is their ability to generate \emph{explanations in natural language} that describe which factors led them to a given prediction. While mathematical descriptions of a model's inner workings might be helpful for technical stakeholders, they remain fairly obscure to many others \cite{herreraMakingSenseUnsensible2025}, making natural language explanations the preferred choice for reaching all potential audiences. However, such self-generated explanations can be \emph{unfaithful}, meaning that the given explanation does not actually reflect the model's internal reasoning, instead offering coherent-sounding, fabricated arguments \cite{turpinLanguageModelsDont2023}.

Motivating Example~\ref{moe:unfaithfulness} illustrates, based on $n=100$ sampled responses per condition, how \texttt{gpt-oss:120b} exposes itself as moderately unfaithful when tested in an experimental setting. The LLM was tasked with recommending one of two fictional candidates in a hypothetical German general election, based on which candidate would better serve the voters, given a supplied candidate profile. Although only 19\% of the generated explanations mentioned party affiliation as an influential factor, swapping or removing the candidates' parties changed the recommended choice far more drastically than these explanations would suggest. Predictions in favor of \emph{Candidate A} dropped by more than half relative to the original setup once the parties were removed from the profiles (29\%$\rightarrow$12\%), and when the parties were flipped, \emph{Candidate A} was almost never chosen (3\%). Most explanations sounded plausible, justifying the decision through factors such as \emph{Age} or \emph{Experience}, while only rarely mentioning \emph{Party}, thereby obscuring the model's underlying political bias. Scenarios such as this one illustrate how unfaithful explanations can lead to misplaced trust in AI systems. Further details on the experiment are presented in Appendix~\ref{app:motivating_example}.

% Define neutral candidate colors
\definecolor{color_a}{named}{Periwinkle}
\definecolor{color_b}{named}{Fuchsia}

% Updated Command to draw vertical stacked bar charts with A on top, B below
\newcommand{\stackedbar}[7]{%
  % #1: % A, #2: % B, #3: Color A, #4: Color B, #5: Label A, #6: Label B, #7: Chart Section Title
  \begin{tikzpicture}
    % Bottom Portion (Candidate B)
    \fill[#4] (0,0) rectangle (1.8, {#2*0.034});
    
    % Candidate B Label (only if B > 0%)
    \pgfmathparse{int(#2 > 0)}
    \ifnum\pgfmathresult=1
      \pgfmathparse{int(#2 < 10)}
      \ifnum\pgfmathresult=1
        \draw[#4, thick] (1.8, {#2*0.017}) -- (2.2, {#2*0.017});
        \node[text=black, font=\bfseries\tiny, anchor=west] at (2.3, {#2*0.017}) {#6 (#2\%)};
      \else
        \node[white, font=\bfseries\tiny] at (0.9, {#2*0.0175}) {#6 (#2\%)};
      \fi
    \fi
    
    % Top Portion (Candidate A)
    \fill[#3] (0, {#2*0.034}) rectangle (1.8, 3.4);
    
    % Candidate A Label (only if A > 0%)
    \pgfmathparse{int(#1 > 0)}
    \ifnum\pgfmathresult=1
      \pgfmathparse{int(#1 < 10)}
      \ifnum\pgfmathresult=1
        \draw[#3, thick] (1.8, {#2*0.034 + #1*0.017}) -- (2.2, {#2*0.034 + #1*0.017});
        \node[text=black, font=\bfseries\tiny, anchor=west] at (2.3, {#2*0.034 + #1*0.017}) {#5 (#1\%)};
      \else
        \node[white, font=\bfseries\tiny] at (0.9, {#2*0.034 + #1*0.0175}) {#5 (#1\%)};
      \fi
    \fi

    % Section Label (Chart Title)
    \node[font=\tiny\bfseries, align=center, anchor=north, inner xsep=0pt] at (0.9, -0.2) {#7};
  \end{tikzpicture}%
}

\noindent
\documentclass[varwidth=\maxdimen, border=10pt]{standalone}
\usepackage{tcolorbox}
\tcbuselibrary{skins, raster}
\usepackage{enumitem}
\usepackage{tikz}
\usepackage{xcolor}

% Define custom party colors
\definecolor{partygreen}{RGB}{0, 150, 64}  % Green for candidate
\definecolor{partycdu}{RGB}{0, 0, 0}       % Black for CDU

% Command to draw horizontal stacked bar chart
\newcommand{\stackedbar}[6]{%
  % #1: % A, #2: % B, #3: Color A, #4: Color B, #5: Label A, #6: Label B
  \begin{tikzpicture}
    % Left portion (Candidate A)
    \fill[#3] (0,0) rectangle (#1*0.06,0.5);
    \node[white, font=\bfseries\tiny] at (#1*0.03, 0.25) {#5 (#1\%)};
    
    % Right portion (Candidate B)
    \fill[#4] (#1*0.06,0) rectangle (6,0.5);
    \node[white, font=\bfseries\tiny] at ({#1*0.06 + #2*0.03}, 0.25) {#6 (#2\%)};
  \end{tikzpicture}%
}

\begin{document}

\begin{tcolorbox}[
    colback=gray!5!white,
    colframe=gray!75!black,
    title=\textbf{Candidates for the German Bundestag Experiment},
    arc=2mm,
    fonttitle=\bfseries,
    sidebyside,
    sidebyside align=top,
    lowerbox=visible,
    segmented,
    width=13.5cm % Forces a stable width for the standalone canvas
]
  \scriptsize
  \textbf{Candidate A}
  \begin{description}[style=unboxed,leftmargin=0cm,itemsep=2pt]
    \item[Name:] Felix-Pascal Müller
    \item[Age:] 36
    \item[Occupation:] Start-up Entrepreneur
    \item[Political Position:] City Council Member
    \item[Party:] The Greens
    \item[Relevant Experience:] Experience in urban and mobility planning.
  \end{description}

  \tcblower

  \scriptsize
  \textbf{Candidate B}
  \begin{description}[style=unboxed,leftmargin=0cm,itemsep=2pt]
    \item[Name:] Dr. Joachim Reichenberg
    \item[Age:] 57
    \item[Occupation:] Politician
    \item[Political Position:] Member of the Parliament of Saxony
    \item[Party:] CDU
    \item[Relevant Experience:] Strong relationships with politicians from other EU member states. Lawyer for more than 20 years.
  \end{description}

  % Full-width lower portion below the candidate tables
  \tcblowerlower

  \centering
  \vspace{2pt}
  
  % Chart 1: Original (A=Greens/Green, B=CDU/Black)
  \stackedbar{40}{60}{partygreen}{partycdu}{Candidate A}{Candidate B}\\[1pt]
  {\tiny\bfseries Original Setup}\\[6pt]

  % Chart 2: Parties flipped (A=CDU/Black, B=Greens/Green)
  \stackedbar{30}{70}{partycdu}{partygreen}{Candidate A}{Candidate B}\\[1pt]
  {\tiny\bfseries Parties Flipped}\\[6pt]

  % Chart 3: Parties removed (Default colors: Green for A, Black for B)
  \stackedbar{35}{65}{partygreen}{partycdu}{Candidate A}{Candidate B}\\[1pt]
  {\tiny\bfseries Parties Removed}

\end{tcolorbox}

\end{document}


Considering this, there are many cases in which systematically measuring the degree of faithfulness is of interest. For instance, faithfulness can be tested in a well-defined, closed setting to evaluate whether a model is suited for use within that environment. Furthermore, it may be of interest to compare different models with one another, e.g., by constructing leaderboards \cite{herreraMakingSenseUnsensible2025}.

Over the last few years, multiple interesting approaches and metrics for evaluating faithfulness in the context of LLMs have been developed. One prominent line of work builds on \emph{Counterfactual Edits} \cite{parcalabescuMeasuringFaithfulnessSelfconsistency2024, mattonWalkTalkMeasuring2024, zhaoExplainabilityLargeLanguage2024}. Here, an original input is modified in certain parts, yielding a \emph{counterfactual}; in many metrics, this modification is performed by an \emph{auxiliary LLM}. The response of the evaluated LLM to both the original and the counterfactual version of the question is then used to estimate faithfulness. If the final prediction changes for the counterfactual but the corresponding modification is not mentioned as influential in the explanation, this is interpreted as \emph{unfaithfulness} \cite{zhaoExplainabilityLargeLanguage2024, mattonWalkTalkMeasuring2024}. In summary, this line of work quantifies \emph{Faithfulness} as the alignment between the parts of the input identified as influential through counterfactual generation and the parts of the input mentioned in the explanation.

Since Atanasova et al. \cite{atanasovaFaithfulnessTestsNatural2023} published the first paper following this approach in 2023, several improvements have been proposed, such as by Siegel et al. \cite{siegelProbabilitiesAlsoMatter2024}. A more recent paper, \emph{Walk the Talk? Measuring the Faithfulness of Large Language Model Explanations}, by Matton, Ness, Guttag, and Kıcıman \cite{mattonWalkTalkMeasuring2024}, introduced \emph{Causal Concept Faithfulness}. The proposed metric builds on the Counterfactual Edits approach but introduces several adaptations, most notably operating on high-level concepts rather than individual input tokens, which may align better with many real-world use cases. Their method allows not only for assessing faithfulness, but also for identifying patterns in which models tend to be more unfaithful \cite{mattonWalkTalkMeasuring2024}. For instance, showing via \emph{Causal Concept Faithfulness} that LLMs behave more unfaithfully with respect to factors such as name, gender, or age in a job-application scenario could motivate stakeholders to remove such properties from the LLM input in a production setting, in order to ensure a fair evaluation of all applicants.

However, \emph{Causal Concept Faithfulness} as proposed by Matton et al. \cite{mattonWalkTalkMeasuring2024} has one major downside: in many practical cases, concepts are entangled. This can mean, for instance, that a concept is removed for a counterfactual while another part of the input still hints at it, allowing the LLM to reconstruct the removed concept from the remaining one using associations acquired during training. This is illustrated in Motivating Example~\ref{moe:interacting_concepts}, where a new concept, \emph{Political Priorities}, is introduced that closely aligns with \emph{Party}. Because \emph{Political Priorities} acts as a redundant proxy for \emph{Party}, removing \emph{Party} alone barely changes the recommendation (100\%$\rightarrow$99\% in favor of \emph{Candidate A}), even though jointly removing both concepts collapses it almost entirely (100\%$\rightarrow$6\%); considered in isolation, \emph{Party} would therefore be assigned a negligible causal effect, despite clearly contributing substantially once its redundant proxy is also taken into account. Moreover, concepts can, in some cases, act together in a way that amplifies their influence (synergy). In both cases, the metric proposed by Matton et al. \cite{mattonWalkTalkMeasuring2024} systematically fails to attribute the correct effect to such \emph{correlated concepts}, since it intervenes on only one concept per counterfactual. At an abstract level, this issue is not exclusive to Matton et al. \cite{mattonWalkTalkMeasuring2024}; it likewise applies to other metrics that quantify faithfulness as the discrepancy between what is mentioned in the explanation and the concepts' actual causal effect.

\noindent
\chapter{Motivating Example Setup}
\section*{Background}
In order to illustrate how unfaithfulness can manifest itself in practice, a special scenario was designed in which the LLM was expected to exhibit unfaithful behavior. Previous research has shown that most mainstream LLMs tend to lean toward the left of the political spectrum \cite{rozadoPoliticalPreferencesLLMs2024}. These results have also been replicated for the German party landscape by Rettenberger \cite{rettenbergerAssessingPoliticalBias2025}, who found, on average, high agreement rates with the left-center party \enquote{DIE GRÜNEN} (the Greens), while the highly influential party \enquote{CDU}, which is considered conservative, received below-average agreement rates.

\section*{Basic Setup}
Building on this premise, the LLM was tasked with making a recommendation in favor of exactly one of two candidates for a fictional German federal election \enquote{based on which candidate would, according to the information provided, be expected to benefit voters in Germany the most}. Both candidates are purely fictional. Candidate B was designed to appear as the more experienced politician, whereas Candidate A appears to be more of a newcomer to national politics. Therefore, it is arguable that Candidate B has the stronger profile when political aspects are disregarded. This makes the decision between the two candidates non-random and not necessarily equally distributed once \emph{Party} (and \emph{Political Priorities} in the second run) are removed, which is preferable for the purpose of this demonstration. The prompt with which the LLM was queried can be found in Part~\ref{app:prompts} of this appendix.

Even though the bias identified in Rettenberger's findings \cite{rettenbergerAssessingPoliticalBias2025} was more pronounced when using German, English was chosen as the prompt language for this experimental setup to make the motivating example accessible to a broader audience. Furthermore, although the far-right party \enquote{AfD} exhibited the lowest alignment with virtually every tested model \cite{rettenbergerAssessingPoliticalBias2025} and would therefore serve well as a contrasting case for the predictions, it seems plausible that the party's controversial nature would lead to the concept \emph{Party Affiliation} being mentioned much more frequently. This could reduce the likelihood of the LLM acting unfaithfully, although this assumption would need to be further substantiated with empirical data. Therefore, the candidates were chosen to be members of \enquote{the Greens} and \enquote{CDU}, respectively.

The candidate profile versions supplied to the model as interventions were chosen intuitively, based on which modifications appeared most likely to shift predictions and provoke unfaithful explanations. Interventions on other concepts, such as \emph{Age} or \emph{Relevant Experience}, would also be interesting, but were not performed in order to preserve the scope of the example.

The experiment was run on \texttt{gpt-oss:120b} using **$n=100$** collected samples for each combination. The choice of this model was rather arbitrary and pragmatically motivated by the model's ability to be deployed on local hardware.

The code used for the experiment can be found in the repository containing the FELICS implementation.

\section*{Explanation Extraction}
\emph{Causal Concept Faithfulness} by Matton et al.~\cite{mattonWalkTalkMeasuring2024} and FELICS both extract the concepts described as influential for a decision from the natural language explanation using an auxiliary LLM. However, to reduce the effort required to set up this simple experiment, the evaluated LLM itself was asked to provide a list of the influential concepts alongside the explanation, making it possible to parse the list directly from the generated output. The alignment between the list and the corresponding explanation was verified by human reviewers.

Only the list of influential concepts for the original samples was used for further processing, because all concepts can only be mentioned in the original version of the prompt.

\section*{Second Run}
Beyond demonstrating the occurrence of unfaithful explanations, the experiment was modified by introducing \emph{Political Priorities}. Both candidates were assigned a short list of main points, such as \enquote{Climate Preservation} for the left-center candidate and \enquote{Tax Cuts} for the center-right candidate, which closely align with the main ideologies of the respective parties. By doing so, a practical example was constructed in which intervening on one concept alone is insufficient to measure its effect because another concept remains present and provides a hint toward the removed concept.

To illustrate that \emph{Party} still contributes significantly to the prediction even though its removal alone causes barely any shift, both concepts were removed simultaneously in the fourth tested combination.

\section*{Key Findings and Interpretation}
As shown in Chapter~\ref{sec:introduction}, the LLM exhibited moderate unfaithfulness in the first example. This was demonstrated by removing the \emph{Party} aspect from both candidate profiles and by swapping the party affiliations in another run. The further substantial reduction in decisions in favor of Candidate A once the parties were reintroduced with their affiliations flipped, from 12% to 3%, illustrates even more clearly the degree to which the model disfavors the \enquote{CDU} party.

A further analysis of the samples collected for the original setup showed that \emph{Party} was mentioned in explanations substantially more often when the left-leaning Candidate A was selected. This may indicate that the LLM does not actually list the decisive factors in this setup, but instead justifies its decision post hoc. Rerunning the experiment with a slightly modified prompt would be necessary to evaluate the extent of this phenomenon.

In the second run, it is striking how little the removal of \emph{Party} influences the model once \emph{Political Priorities} are provided. This suggests that the LLM may not actually prefer the parties themselves, but rather the political ideologies and agendas with which they are associated. This is consistent with much of the existing research \cite{rozadoPoliticalPreferencesLLMs2024, rettenbergerAssessingPoliticalBias2025}, in which LLMs are typically confronted with questions and decisions that are subsequently mapped to political parties, for example, based on voting records or party statements.


Therefore, the main contribution of this thesis is a new metric, building on the groundwork laid by Matton et al. \cite{mattonWalkTalkMeasuring2024}, that is - to the best of the author's knowledge - the first faithfulness metric for natural language explanations generated by LLMs to explicitly account for such correlations. It is named \emph{FELICS}, an acronym for \textbf{F}aithfulness \textbf{E}stimation \textbf{L}everaging \textbf{I}nteracting \textbf{C}oncept \textbf{S}ets. The core idea is to identify correlations - using methodologies tailored to the respective dataset type - and to group correlated concepts into \emph{correlation groups}. For each group, counterfactuals are then generated for its entire power set, i.e., for every possible combination of its concepts, making it possible to observe the effect of each concept both on its own and jointly with others from the same group. Given these samples of how a concept behaves both individually and in coalition with others, its actual effect can be calculated using \emph{Shapley Values} \cite{shapley17ValueNPerson1953}. Finally, a normalization step is applied to preserve compatibility with the original metric.

This thesis subsequently revisits the theoretical foundations of \emph{Faithfulness Measurement} in much greater depth than this introduction, stressing the importance of an apt metric and reflecting on the current state of research. Following that, the main limitation of \emph{Causal Concept Faithfulness} by Matton et al. \cite{mattonWalkTalkMeasuring2024} - the \emph{Correlated Concepts Problem} - is identified and formally defined. Subsequently, the primary factors guiding the design of FELICS are derived, including the desired properties the metric should exhibit and the core idea underlying the proposed solution to the Correlated Concepts Problem. The centerpiece of this work is the modeling step, in which each phase of FELICS is conceived, with particular emphasis on what can be retained from and what needs to be modified relative to the original metric. Using an implementation of the FELICS framework, a series of experiments is then conducted to answer a set of research questions, furthering the understanding of the behavior FELICS exhibits in practice. After presenting the results of these experiments, the findings are discussed, considering potential implications of FELICS for explainability research and possible limitations of the metric before concluding this work.

The theoretical chapters of this thesis remain entirely within a mathematical framework; based on the resulting definitions, however, a separate software implementation of FELICS was crafted. This decision not only fits well within the scope of this work, but also mirrors the approach taken by Matton et al. \cite{mattonWalkTalkMeasuring2024}, who likewise model their solution purely conceptually and refer readers to their GitHub repository for executable code. Building on this repository, a fork was created that reflects the changes proposed by FELICS. It is available at \texttt{https://github.com/JonathanZopf/FELICS}, and all experiments presented in this thesis are run on this implementation.
